% ============================================================
% Physics-Constrained Multimodal Lunar Geochemical Retrieval
% IIRS + CLASS-XRF Research Methodology
% Compile with: pdflatex -> bibtex -> pdflatex -> pdflatex
% ============================================================

\documentclass[11pt,a4paper]{article}

\usepackage[margin=1in]{geometry}
\usepackage{amsmath,amssymb,amsfonts,bm}
\usepackage{booktabs}
\usepackage{array}
\usepackage{longtable}
\usepackage{graphicx}
\usepackage{xcolor}
\usepackage{hyperref}
\usepackage{enumitem}
\usepackage{siunitx}
\usepackage{microtype}
\usepackage{float}
\usepackage{listings}
\usepackage{caption}
\usepackage{cite}

\hypersetup{
    colorlinks=true,
    linkcolor=blue!60!black,
    citecolor=blue!60!black,
    urlcolor=blue!70!black
}

\lstset{
    basicstyle=\ttfamily\small,
    breaklines=true,
    frame=single,
    columns=fullflexible
}

\title{\textbf{Physics-Constrained Multimodal Retrieval of Lunar Elemental Ratios}\\
\large A Research-Grade Methodology for Chandrayaan-2 IIRS and CLASS-XRF Data Fusion}

\author{}
\date{\today}

\begin{document}

\maketitle

\begin{abstract}
Near-infrared (NIR) spectroscopy and X-ray fluorescence (XRF) provide complementary measurements of lunar composition. NIR observations from instruments such as the Chandrayaan-2 Imaging Infrared Spectrometer (IIRS) primarily constrain mineralogy, mineral chemistry, grain size, space-weathering state, and vibrational/electronic absorption features, whereas the Chandrayaan-2 Large Area Soft X-ray Spectrometer (CLASS) directly measures elemental X-ray fluorescence lines for major elements under suitable solar illumination. Consequently, elemental ratios such as Mg/Si, Al/Si, Ca/Si, and Fe/Si should not be interpreted as native NIR observables.

This methodology formulates the problem as a multimodal inverse problem. IIRS spectra are converted into calibrated reflectance and then analyzed using classical absorption parameters, nonlinear radiative-transfer spectral unmixing, and physics-aware machine-learning representations. CLASS spectra are independently processed through solar-state screening, particle-background subtraction, instrumental response correction, and XRF line fitting. The two data products are subsequently harmonized in spatial scale and fused using Bayesian or physics-constrained machine-learning inference. Grain size, space weathering, mineral mixing, photometric geometry, thermal emission, calibration uncertainty, and model discrepancy are explicitly treated as nuisance or latent variables.

The final data product consists of elemental ratios with uncertainty distributions, detection limits, quality flags, provenance, and validation pedigree. NIR-only products are restricted to mineral-chemical quantities such as Mg\#, Fo\#, En\#, Wo\#, modal mineral fractions, and FeO proxies. Absolute elemental ratios are reported only after calibration or fusion with CLASS-XRF and independent geochemical constraints.
\end{abstract}

\tableofcontents
\newpage

% ============================================================
\section{Scientific Scope and Measurement Philosophy}
% ============================================================

\subsection{Problem Definition}

The objective is to retrieve lunar elemental ratios from complementary hyperspectral and XRF measurements while preserving the distinction between directly measured quantities and model-dependent proxies.

For the NIR measurement, the forward model is written as

\begin{equation}
\mathbf{y}_{\mathrm{NIR}}
=
\mathcal{F}_{\mathrm{NIR}}
\left(
\mathbf{m},
\mathbf{g},
\mathbf{w},
\mathbf{p},
T
\right)
+
\boldsymbol{\epsilon}_{\mathrm{NIR}},
\end{equation}

where $\mathbf{m}$ represents modal mineral abundances, $\mathbf{g}$ mineral chemistry, $\mathbf{w}$ space-weathering state, $\mathbf{p}$ photometric geometry, $T$ surface temperature, and $\boldsymbol{\epsilon}$ measurement noise.

For CLASS-XRF,

\begin{equation}
\mathbf{y}_{\mathrm{XRF}}
=
\mathcal{F}_{\mathrm{XRF}}
\left(
\mathbf{c},
\mathbf{G},
\mathbf{S},
\mathbf{B}
\right)
+
\boldsymbol{\epsilon}_{\mathrm{XRF}},
\end{equation}

where $\mathbf{c}$ is elemental abundance, $\mathbf{G}$ observation geometry, $\mathbf{S}$ the incident solar spectrum/state, and $\mathbf{B}$ the instrumental and particle background.

The desired posterior is therefore

\begin{equation}
P
\left(
\mathbf{m},
\mathbf{g},
\mathbf{c},
\mathbf{w}
\mid
\mathbf{y}_{\mathrm{NIR}},
\mathbf{y}_{\mathrm{XRF}}
\right).
\end{equation}

This formulation is preferable to a simple regression from NIR band parameters to elemental ratios because it explicitly represents the degeneracies inherent in reflectance spectroscopy.

\subsection{Measurement Complementarity}

NIR absorptions principally arise from crystal-field transitions, charge-transfer processes, and molecular vibrational overtones/combination bands. They therefore constrain mineralogical state rather than directly counting atoms.

CLASS-XRF, by contrast, provides direct information on characteristic X-ray fluorescence lines. The methodological distinction is:

\begin{equation}
\boxed{
\text{NIR} \rightarrow \text{mineralogy and mineral chemistry}
}
\end{equation}

\begin{equation}
\boxed{
\text{CLASS-XRF} \rightarrow \text{elemental abundance}
}
\end{equation}

The research problem is consequently a \emph{data-fusion inverse problem}, not merely an NIR elemental-ratio regression.

% ============================================================
\section{Data Sources and Instrument-Specific Products}
% ============================================================

\subsection{Chandrayaan-2 IIRS}

IIRS data should be treated as a calibrated hyperspectral data cube rather than a collection of isolated band measurements. The operational chain should preserve the complete spectrum wherever possible.

The useful spectral regimes are divided into:

\begin{enumerate}[label=\alph*.]
    \item approximately $0.8$--$2.5~\mu$m: primary mineralogical VNIR/SWIR regime;
    \item approximately $2.5$--$3.3~\mu$m: strongly coupled mineralogical, hydration, and thermal regime;
    \item approximately $3.3$--$5.0~\mu$m: increasingly important thermal/emissive regime.
\end{enumerate}

Recent IIRS Level-2 processing work emphasizes radiometric correction, thermal correction, apparent reflectance, standardized geometry, empirical-line correction, and validation against independent lunar observations. These steps should be reflected explicitly in the pipeline.

\subsection{CLASS-XRF}

CLASS measurements should be processed as event/spectrum data with explicit solar-state and particle-background quality control. Elemental inference should be based on physically modeled characteristic lines rather than simple raw count ratios.

Relevant characteristic lines include Mg, Al, Si, Ca, and Fe fluorescence features. The final XRF elemental retrieval should incorporate:

\begin{itemize}
    \item detector response;
    \item energy calibration;
    \item line-spread function;
    \item solar spectral state;
    \item continuum background;
    \item particle-induced background;
    \item pile-up and instrumental artifacts;
    \item fluorescence line fitting;
    \item uncertainty in fundamental-parameter conversion.
\end{itemize}

% ============================================================
\section{IIRS Level-1 to Level-2 Processing}
% ============================================================

\subsection{Detector and Radiometric Calibration}

Starting from detector digital numbers,

\begin{equation}
DN_1
=
DN_{\mathrm{raw}}
-
D(T_{\mathrm{det}})\,t_{\mathrm{exp}}
-
B_{\mathrm{bias}},
\end{equation}

followed by flat-field correction,

\begin{equation}
DN_2 =
\frac{DN_1}{F_{\mathrm{flat}}(x,\lambda)}.
\end{equation}

Bad pixels, dropouts, cosmic-ray spikes, nonlinear detector response, and saturation should be flagged rather than silently interpolated.

Radiance is obtained through an instrument calibration function,

\begin{equation}
L_{\mathrm{rad}}(\lambda)
=
DN_2 G(\lambda),
\end{equation}

where $G(\lambda)$ incorporates the relevant calibration coefficients.

\subsection{Wavelength Calibration}

The wavelength solution may initially be represented as

\begin{equation}
\lambda_0(p)
=
a_0+a_1p+a_2p^2+\cdots,
\end{equation}

with detector temperature, optical flexure, smile, and spatially dependent wavelength offsets included in the operational calibration.

The corrected spectrum may be resampled onto a common wavelength grid using a flux-preserving interpolation scheme.

The residual wavelength error should be propagated:

\begin{equation}
\sigma_{\lambda,\mathrm{total}}^2
=
\sigma_{\lambda,\mathrm{fit}}^2
+
\sigma_{\lambda,\mathrm{thermal}}^2
+
\sigma_{\lambda,\mathrm{cal}}^2.
\end{equation}

\subsection{Thermal Emission Removal}

For wavelengths where thermal emission becomes significant,

\begin{equation}
L_{\mathrm{obs}}
=
L_{\mathrm{reflected}}
+
\epsilon(\lambda)B_\lambda(T),
\end{equation}

where $B_\lambda(T)$ is the Planck function.

Thermal removal should be iterative because errors in $T$ and emissivity propagate strongly into the $2.5$--$5~\mu$m region.

The retrieved temperature should therefore be treated as an inferred nuisance parameter with uncertainty rather than as an exact external quantity.

\subsection{Photometric Correction}

The reflectance should be normalized to a common geometry using an appropriate photometric model such as Hapke, Lommel--Seeliger, Akimov, or an empirically validated lunar disk function.

\begin{equation}
R_{\mathrm{std}}(\lambda)
=
\mathcal{P}
\left[
R_{\mathrm{obs}}(\lambda);
i,e,g
\rightarrow
i_0,e_0,g_0
\right].
\end{equation}

The photometric correction must not be assumed to remove space-weathering effects.

% ============================================================
\section{Spectral Feature Extraction}
% ============================================================

\subsection{Continuum Removal}

For an absorption feature,

\begin{equation}
R_n(\lambda)
=
\frac{R(\lambda)}{R_c(\lambda)},
\end{equation}

where $R_c$ is a locally estimated continuum.

Band depth is

\begin{equation}
BD(\lambda_c)
=
1-R_n(\lambda_c).
\end{equation}

Equivalent width is

\begin{equation}
EW
=
\int_{\lambda_1}^{\lambda_2}
\left[1-R_n(\lambda)\right]d\lambda.
\end{equation}

Band area may be used to calculate

\begin{equation}
BAR
=
\frac{BA_{II}}{BA_I}.
\end{equation}

Continuum methods should include convex-hull, polynomial, and physically motivated local baselines. The difference between methods is itself a systematic uncertainty.

\subsection{Classical Mineralogical Parameters}

The principal parameters include

\[
\lambda_{BI},\quad
\lambda_{BII},\quad
BD_I,\quad
BD_{II},\quad
BA_I,\quad
BA_{II},\quad
BAR,
\]

along with albedo, continuum slope, spectral curvature, and hydration-region parameters.

These quantities should be retained even when a full-spectrum machine-learning representation is subsequently used.

\subsection{Full-Spectrum Representation}

The complete spectrum should also be represented using methods such as:

\begin{itemize}
    \item PCA/MNF;
    \item wavelet representations;
    \item autoencoder embeddings;
    \item physically constrained latent representations;
    \item nonlinear radiative-transfer parameters.
\end{itemize}

This avoids the information loss associated with reducing every spectrum to a small number of band parameters.

% ============================================================
\section{Nonlinear Spectral Unmixing}
% ============================================================

\subsection{Motivation}

Linear spectral mixing is a useful baseline, but lunar regolith is an intimate particulate mixture for which nonlinear scattering can be important. The final inversion should therefore compare linear mixing with nonlinear radiative-transfer models.

A generic physical mixture model is

\begin{equation}
R_{\mathrm{model}}(\lambda)
=
\mathcal{H}
\left(
f_1,\ldots,f_N,
D,
\theta,
w
\right),
\end{equation}

where $f_j$ are modal abundances, $D$ represents grain-size parameters, $\theta$ represents photometric/radiative-transfer parameters, and $w$ represents optical maturity or space-weathering state.

\subsection{Hapke/Shkuratov Inversion}

The inversion minimizes

\begin{equation}
\chi^2
=
\sum_{\lambda}
\frac{
\left[
R_{\mathrm{obs}}(\lambda)
-
R_{\mathrm{model}}(\lambda)
\right]^2
}{
\sigma_R^2(\lambda)
}.
\end{equation}

Constraints include

\begin{equation}
f_j\geq0,
\end{equation}

and

\begin{equation}
\sum_j f_j=1.
\end{equation}

The solution should be reported as a posterior distribution rather than as a single deterministic vector whenever computationally feasible.

\subsection{Machine-Learning Surrogate}

Recent research combining nonlinear Shkuratov modeling with machine learning provides a useful framework for scaling physically based inversion to large hyperspectral datasets.

A surrogate model may be defined as

\begin{equation}
\mathbf{z}
=
NN_\theta
\left[
R(\lambda)
\right],
\end{equation}

where $\mathbf{z}$ contains mineral fractions and mineral-chemical parameters.

The ML model should be trained on synthetic spectra generated from the physical forward model, supplemented by laboratory and flight observations.

% ============================================================
\section{Explicit Treatment of Space Weathering and Grain Size}
% ============================================================

Space weathering should be treated as a latent variable rather than simply removed as a preprocessing nuisance.

The reflectance model is therefore

\begin{equation}
R(\lambda)
=
F
\left(
\text{mineralogy},
\text{chemistry},
\text{grain size},
\text{maturity},
\text{geometry}
\right).
\end{equation}

A useful parameter vector is

\begin{equation}
\boldsymbol{\theta}
=
\left[
f_{\mathrm{ol}},
f_{\mathrm{opx}},
f_{\mathrm{cpx}},
f_{\mathrm{plg}},
Fo\#,
En\#,
Wo\#,
D,
\tau_{\mathrm{weathering}}
\right].
\end{equation}

Inference should marginalize over $D$ and $\tau_{\mathrm{weathering}}$ rather than assuming they are known.

This is essential because grain size, optical maturity, albedo, continuum slope, and absorption strength can produce degeneracies that otherwise appear as compositional trends.

% ============================================================
\section{NIR-Derived Mineral Chemistry and Proxy Products}
% ============================================================

\subsection{Olivine}

Olivine composition should primarily be expressed using Fo\# or Fa\# rather than directly claiming a measured Mg/Si ratio.

\begin{equation}
Fo\#
=
100
\frac{Mg}{Mg+Fe}.
\end{equation}

The conversion from spectral band position to Fo\# should be treated as a calibration relation with uncertainty, not a universal exact equation.

\subsection{Pyroxene}

For pyroxene, the inversion should estimate En, Fs, and Wo components where the calibration permits:

\begin{equation}
En+Fs+Wo=100.
\end{equation}

Band I and Band II centers, areas, widths, and asymmetries should be used jointly.

\subsection{Feldspathic Material}

Plagioclase is difficult to quantify from NIR alone because relatively Fe-poor feldspar can have weak absorptions.

Consequently, NIR should report:

\begin{itemize}
    \item feldspathic probability;
    \item modal plagioclase estimate with uncertainty;
    \item anorthite-related classification where justified;
    \item confidence/upper-limit flags.
\end{itemize}

Al/Si should not be treated as a direct NIR measurement.

% ============================================================
\section{Elemental-Ratio Feasibility}
% ============================================================

\begin{table}[H]
\centering
\caption{Recommended interpretation of NIR and CLASS observables.}
\begin{tabular}{p{0.14\textwidth}p{0.28\textwidth}p{0.22\textwidth}p{0.23\textwidth}}
\toprule
Ratio & NIR information & NIR interpretation & CLASS interpretation \\
\midrule
Mg/Si & Olivine/pyroxene chemistry, Mg\#, Fo\#, En\# & Model-dependent proxy & Direct XRF elemental constraint \\
Fe/Si & Band strengths, Fe-bearing mineral abundance, FeO proxy & Strongly affected by maturity and grain size & Direct XRF elemental constraint \\
Ca/Si & Pyroxene Wo\#, HCP/LCP discrimination & Mineralogical proxy; plagioclase correction required & Direct XRF elemental constraint \\
Al/Si & Feldspathic probability, weak Fe-bearing feldspar feature & Low-precision indirect proxy & Direct XRF elemental constraint \\
\bottomrule
\end{tabular}
\end{table}

NIR-only elemental ratios should therefore be labeled explicitly as \emph{model-dependent reconstructed ratios}. True elemental ratios should be produced only after external calibration or multimodal fusion.

% ============================================================
\section{CLASS-XRF Processing}
% ============================================================

\subsection{Solar-State Screening}

CLASS XRF intensity is strongly dependent on incident solar X-ray flux. Observations should therefore be classified according to solar state using contemporaneous solar-monitor information.

Only observations meeting the required solar-flux and particle-background quality criteria should contribute to quantitative elemental inversion.

\subsection{Energy Calibration}

The detector energy scale is represented by

\begin{equation}
E(p)
=
b_0+b_1p+b_2p^2+\cdots,
\end{equation}

where $p$ is detector channel.

Energy drift should be tracked with time and instrument state.

\subsection{Line Fitting}

The observed spectrum can be represented as

\begin{equation}
C(E)
=
B(E)
+
\sum_k
A_k L_k(E)
+
P(E),
\end{equation}

where $B(E)$ is the continuum/background, $L_k(E)$ are instrument-broadened fluorescence lines, and $P(E)$ represents particle-induced components.

The fitting process should return both line fluxes and their covariance matrix.

\subsection{Fundamental-Parameter Inversion}

Elemental abundance is inferred from the measured line intensities using a physical XRF response model,

\begin{equation}
\mathbf{F}
=
\mathcal{G}
\left(
\mathbf{c},
\mathbf{S},
\mathbf{G},
\mathbf{R}_{det}
\right).
\end{equation}

The inverse solution is

\begin{equation}
P(\mathbf{c}\mid\mathbf{F})
\propto
P(\mathbf{F}\mid\mathbf{c})
P(\mathbf{c}).
\end{equation}

% ============================================================
\section{Spatial Harmonization of IIRS and CLASS}
% ============================================================

IIRS and CLASS operate at substantially different spatial scales. Direct pixel-to-pixel comparison is therefore invalid.

For a CLASS footprint $k$, aggregate valid IIRS measurements using uncertainty and quality weights:

\begin{equation}
\bar{P}_{\mathrm{IIRS},k}
=
\frac{
\sum_i w_i P_i
}{
\sum_iw_i
}.
\end{equation}

Weights should incorporate:

\begin{equation}
w_i
\propto
\frac{
q_i
}{
\sigma_i^2
},
\end{equation}

where $q_i$ is a quality/homogeneity factor.

The aggregation should also test lithological homogeneity. A CLASS footprint containing multiple lithologies should not be treated as a single pure calibration target.

The covariance generated by spatial aggregation must be retained.

% ============================================================
\section{Multimodal Fusion}
% ============================================================

\subsection{Bayesian Fusion}

The central fusion equation is

\begin{equation}
P(\mathbf{c}\mid
\mathbf{y}_{NIR},
\mathbf{y}_{XRF})
\propto
P(\mathbf{y}_{NIR}\mid\mathbf{c})
P(\mathbf{y}_{XRF}\mid\mathbf{c})
P(\mathbf{c}).
\end{equation}

The NIR likelihood constrains mineralogical and chemical consistency, while the XRF likelihood anchors elemental abundance.

\subsection{Physics-Constrained Neural Fusion}

A neural architecture may define

\begin{equation}
\mathbf{z}_{NIR}
=
f_\theta(R_{IIRS}),
\end{equation}

\begin{equation}
\mathbf{z}_{XRF}
=
g_\phi(F_{Mg},F_{Al},F_{Si},F_{Ca},F_{Fe}),
\end{equation}

and

\begin{equation}
\mathbf{z}
=
[
\mathbf{z}_{NIR},
\mathbf{z}_{XRF},
\mathbf{G}
].
\end{equation}

The final elemental prediction is

\begin{equation}
\hat{\mathbf{c}}
=
h_\psi(\mathbf{z}).
\end{equation}

The loss function should contain observational, physical, stoichiometric, and spatial terms:

\begin{equation}
\mathcal{L}
=
\mathcal{L}_{spectral}
+
\lambda_1\mathcal{L}_{XRF}
+
\lambda_2\mathcal{L}_{stoich}
+
\lambda_3\mathcal{L}_{physics}
+
\lambda_4\mathcal{L}_{spatial}.
\end{equation}

\subsection{Physical Constraints}

The network should satisfy

\begin{equation}
f_i\geq0,
\qquad
\sum_i f_i=1,
\end{equation}

and elemental abundances must remain non-negative.

Stoichiometric reconstruction is represented by

\begin{equation}
\mathbf{c}_{NIR}
=
\mathbf{S}\mathbf{f},
\end{equation}

where $\mathbf{S}$ is the mineral stoichiometry matrix.

The predicted composition should also satisfy geochemical closure within the uncertainty associated with unmodeled volatile/light-element components.

% ============================================================
\section{Uncertainty Quantification}
% ============================================================

\subsection{Measurement Uncertainty}

Detector and photon noise should be propagated through the complete calibration chain.

\begin{equation}
\sigma_{\mathrm{shot}}
=
\frac{
\sqrt{
N_e+N_{\mathrm{dark}}+RN^2
}
}{
N_e
}.
\end{equation}

\subsection{Band-Depth Uncertainty}

For

\begin{equation}
BD=1-\frac{R_b}{R_c},
\end{equation}

the variance is

\begin{equation}
\sigma_{BD}^2
=
\frac{\sigma_{R_b}^2}{R_c^2}
+
\frac{R_b^2\sigma_{R_c}^2}{R_c^4}
-
2\frac{R_b}{R_c^3}
\operatorname{Cov}(R_b,R_c).
\end{equation}

Correlated calibration errors must therefore not be treated as independent channel noise.

\subsection{Monte Carlo Retrieval}

For final products, use correlated perturbations:

\begin{lstlisting}
for k = 1 ... N:
    R_k       = R + correlated_noise(C_R)
    lambda_k  = lambda + noise(sigma_lambda)
    anchors_k = perturb_continuum_anchors()
    Rc_k      = fit_continuum(R_k, anchors_k)
    params_k  = extract_features(R_k, Rc_k)
    chem_k    = nonlinear_inversion(params_k, R_k)

report median and 16/84 percentiles
\end{lstlisting}

At least $N=500$ realizations are recommended for routine retrieval and approximately $N=1000$ for high-value science maps.

\subsection{Three-Level Uncertainty Budget}

Separate:

\begin{enumerate}
    \item measurement uncertainty;
    \item retrieval uncertainty;
    \item geological/model uncertainty.
\end{enumerate}

This prevents model inadequacy from being incorrectly represented as detector noise.

% ============================================================
\section{Detection Limits and Quality Flags}
% ============================================================

A conservative detection threshold is

\begin{equation}
BD_{\mathrm{lim}}
=
3\sigma_{BD,\mathrm{MC}}.
\end{equation}

If

\begin{equation}
BD < BD_{\mathrm{lim}},
\end{equation}

the product should report an upper limit rather than a numerical composition.

Recommended quality flags are:

\begin{itemize}
    \item \texttt{GOOD};
    \item \texttt{SNR\_LIMITED};
    \item \texttt{THERMAL\_CONTAMINATED};
    \item \texttt{SATURATED};
    \item \texttt{MIXTURE\_AMBIGUOUS};
    \item \texttt{DETECTION\_LIMIT};
    \item \texttt{OUTSIDE\_CALIBRATION}.
\end{itemize}

% ============================================================
\section{Laboratory and External Validation}
% ============================================================

\subsection{Laboratory Spectral Libraries}

Candidate libraries include RELAB, USGS spectral libraries, and HOSERLab measurements.

The laboratory spectra should be transformed into flight-like observations by:

\begin{enumerate}
    \item spectral convolution to IIRS resolution;
    \item wavelength-grid resampling;
    \item noise injection;
    \item grain-size variation;
    \item space-weathering variation;
    \item thermal and photometric perturbation where appropriate.
\end{enumerate}

The pipeline should then recover the known mineral chemistry within the declared uncertainty.

\subsection{Cross-Instrument Validation}

The validation hierarchy should include:

\begin{enumerate}
    \item IIRS internal calibration and processing validation;
    \item M$^3$ spectral overlap;
    \item laboratory mineral spectra;
    \item CLASS-XRF elemental products;
    \item lunar sample geochemistry;
    \item orbital gamma-ray measurements for independent regional-scale validation;
    \item independent mineralogical maps.
\end{enumerate}

\subsection{Blind Spatial Validation}

Random pixel splitting should be avoided because neighboring lunar pixels are spatially correlated.

Instead, use geographically separated training, validation, and test regions:

\begin{lstlisting}
TRAIN
  - geological region A
  - geological region B
  - geological region C

VALIDATION
  - independent region D

TEST
  - geographically separated region E
\end{lstlisting}

The test region must not participate in calibration or hyperparameter selection.

% ============================================================
\section{Model Comparison and Ablation Study}
% ============================================================

The research should benchmark increasingly sophisticated models:

\begin{table}[H]
\centering
\caption{Recommended model hierarchy.}
\begin{tabular}{p{0.35\textwidth}p{0.50\textwidth}}
\toprule
Model & Scientific purpose \\
\midrule
Band-center regression & Classical spectroscopy baseline \\
PLSR & Linear statistical baseline \\
Random Forest / XGBoost & Nonlinear statistical baseline \\
Linear spectral unmixing & Mineralogical baseline \\
Hapke/Shkuratov inversion & Physics-based baseline \\
Shkuratov + ML surrogate & Scalable nonlinear baseline \\
Physics-constrained neural inversion & Proposed advanced model \\
IIRS + CLASS multimodal fusion & Final research model \\
\bottomrule
\end{tabular}
\end{table}

Ablation studies should remove one component at a time:

\begin{itemize}
    \item no thermal correction;
    \item no weathering parameter;
    \item no grain-size parameter;
    \item linear instead of nonlinear mixing;
    \item no CLASS information;
    \item no spatial harmonization;
    \item no physics constraints.
\end{itemize}

This demonstrates which methodological components actually improve retrieval accuracy.

% ============================================================
\section{Closure and Sanity Tests}
% ============================================================

\subsection{Synthetic Injection-Recovery}

Known absorptions should be injected into representative spectra across SNR values of approximately 10--200.

The retrieval should satisfy:

\begin{equation}
|\Delta\lambda_c|
<
\sigma_{\lambda}
\end{equation}

and, for statistically significant bands,

\begin{equation}
\frac{
|\Delta BD|
}{
BD
}
<
10\%.
\end{equation}

\subsection{Continuum Sensitivity}

Reprocess scenes with:

\begin{itemize}
    \item convex hull;
    \item second-order polynomial;
    \item perturbed anchor points.
\end{itemize}

If the composition changes by more than the reported uncertainty, the method uncertainty must be inflated.

\subsection{Spatial Closure}

Adjacent orbital observations should agree within combined uncertainty. Track-boundary steps indicate residual calibration or photometric errors.

\subsection{Geochemical Closure}

Reconstructed oxide abundances should satisfy an approximate closure condition,

\begin{equation}
\sum_i X_i
=
100\pm\delta~\mathrm{wt.\%},
\end{equation}

where $\delta$ explicitly accounts for excluded volatile/light-element components and model uncertainty.

% ============================================================
\section{Final Data Products}
% ============================================================

The proposed processing levels are:

\begin{description}
    \item[L1:] calibrated radiance / $I/F$ with detector, wavelength, and calibration uncertainties;
    \item[L2:] thermal- and photometrically corrected reflectance plus spectral parameters;
    \item[L3:] mineralogical retrieval including modal fractions, Mg\#, Fo\#, En\#, Wo\#, FeO proxy, grain size, and maturity posterior;
    \item[L4:] CLASS-calibrated elemental ratios Mg/Si, Al/Si, Ca/Si, and Fe/Si;
    \item[L5:] multimodal geochemical maps with posterior uncertainty, detection limits, quality flags, and provenance.
\end{description}

Every L4/L5 elemental-ratio pixel should contain:

\begin{itemize}
    \item posterior median;
    \item $1\sigma$ or 68\% credible interval;
    \item 95\% credible interval;
    \item detection limit;
    \item quality flag;
    \item calibration version;
    \item number of Monte Carlo samples;
    \item continuum method;
    \item spectral library version;
    \item CLASS observation/solar-state metadata;
    \item spatial aggregation footprint;
    \item validation status.
\end{itemize}

% ============================================================
\section{Research Hypothesis and Expected Contribution}
% ============================================================

The central research hypothesis is:

\begin{quote}
A physics-constrained multimodal inversion combining high-spatial-resolution lunar NIR spectroscopy with lower-spatial-resolution but compositionally direct XRF measurements can retrieve elemental ratios with lower systematic bias and better quantified uncertainty than either NIR-only empirical regression or CLASS-only interpolation.
\end{quote}

The principal scientific contribution is therefore not simply an improved regression equation. It is a framework in which:

\begin{equation}
\boxed{
\text{IIRS mineralogical information}
+
\text{CLASS elemental information}
+
\text{physical forward modeling}
+
\text{uncertainty-aware fusion}
}
\end{equation}

produces spatially resolved and physically interpretable lunar geochemical products.

% ============================================================
\section{Implementation Blueprint}
% ============================================================

\begin{lstlisting}
INPUT
  IIRS Level-1 radiance
  CLASS XRF spectra/events
  SPICE geometry
  DEM/topography
  solar-state data
  laboratory spectral libraries

IIRS PIPELINE
  detector calibration
  wavelength calibration
  bad-pixel masking
  thermal correction
  photometric correction
  standardized reflectance

CLASS PIPELINE
  energy calibration
  solar-state screening
  particle-background subtraction
  detector-response correction
  XRF line fitting
  elemental inversion

SPECTRAL PIPELINE
  continuum removal
  band parameters
  full-spectrum embedding
  nonlinear spectral unmixing
  grain-size retrieval
  space-weathering retrieval
  mineral chemistry

SPATIAL PIPELINE
  co-registration
  IIRS -> CLASS footprint aggregation
  homogeneity testing
  covariance propagation

FUSION
  Bayesian inversion OR
  physics-constrained neural model

OUTPUT
  mineral chemistry
  Mg/Si
  Al/Si
  Ca/Si
  Fe/Si
  uncertainty posterior
  detection limits
  quality flags
  provenance

VALIDATION
  laboratory recovery
  synthetic injection-recovery
  M3/IIRS comparison
  CLASS comparison
  lunar sample comparison
  independent regional blind test
  geochemical closure
\end{lstlisting}

% ============================================================
\section{Selected Research Literature}
% ============================================================

\begin{thebibliography}{99}

\bibitem{IIRS2025}
Recent IIRS Level-2 processing work,
\emph{Meteoritics \& Planetary Science},
2025.
Chandrayaan-2 IIRS radiometric/thermal correction, apparent reflectance,
standardized geometry, and validation.

\bibitem{ShkuratovML2024}
Recent study combining the Shkuratov nonlinear mixing model and machine
learning for large-scale lunar M$^3$ compositional mapping,
2024.
The approach explicitly addresses nonlinear mixing, grain size, and optical
maturity.

\bibitem{CLASS2023}
Recent CLASS elemental mapping study,
\emph{Icarus},
2023.
Global lunar elemental mapping using Chandrayaan-2 CLASS X-ray fluorescence
measurements, including Mg, Al, Si, Ca, and Fe.

\bibitem{CLASS2025}
Recent CLASS XRF elemental-ratio study,
2025.
Derivation of lunar O/Si, Mg/Si, Al/Si, Mg/Al, Ca/Si, and Fe/Si from CLASS
fluorescence measurements.

\bibitem{Moriarty2016}
Moriarty, D. P., Pieters, C. M., Isaacson, P. J., and others.
Pyroxene composition and spectral band-center relationships in lunar
materials.
\emph{Meteoritics \& Planetary Science},
2016.

\bibitem{Lucey2000}
Lucey, P. G., Blewett, D. T., and Jolliff, B. L.
Lunar iron and titanium abundance algorithms based on final Clementine
imaging data.
\emph{Journal of Geophysical Research},
2000.

\bibitem{Hapke2012}
Hapke, B.
\emph{Theory of Reflectance and Emittance Spectroscopy}.
Cambridge University Press,
2012.

\bibitem{Adams1974}
Adams, J. B.
Visible and near-infrared diffuse reflectance spectra of pyroxenes as
applied to remote sensing of solid objects in the solar system.
\emph{Journal of Geophysical Research},
1974.

\bibitem{Cloutis1990}
Cloutis, E. A. and Gaffey, M. J.
Pyroxene spectroscopy and the remote sensing of planetary surfaces.
Relevant laboratory and compositional calibration studies.

\bibitem{RELAB}
RELAB Spectral Database.
Brown University.
Laboratory reflectance measurements for planetary spectroscopy.

\bibitem{USGS}
USGS Digital Spectral Library.
United States Geological Survey.
Laboratory mineral spectral reference data.

\end{thebibliography}

% ============================================================
\section*{Methodological Status}
% ============================================================


\end{document}
